\chapter*{この本について}
\addcontentsline{toc}{chapter}{この本について}
\label{ch:about}

この本は、2026 年 10 月 9 日から 21 日までに行った研究「LLM の見通しを Black--Litterman モデル（BLM）に入れる」のすべての記録を、1 冊にまとめたものである。それまでに書いた文書（総合報告書 1・2、実験 1〜11 と補足実験のレポート、2 つの解説、方針と文献調査、設計書、事前登録、結果の記録）をすべて取り込み、金融工学と統計と言語モデルの背景知識の章を新しく書き足して、最初から順に読めるように並べ直した。

\section*{この研究を一言で}

\begin{keybox}[問いと答え]
\textbf{問い}：大規模言語モデル（LLM）に株式の見通しを出させ、それを Black--Litterman モデルのビューとして入れると、時価総額加重の指数（\SP{}）を上回るポートフォリオが作れるか。

\textbf{答え}：作れなかった。ただし理由は段階ごとに違う。
\begin{enumerate}
\item 元論文のやり方（銘柄ごとに翌月のリターンを数字で聞く）では、LLM は入力した直近の値動きをなぞるだけで、固有の予測力は検出できなかった（\textbf{総合報告書 1、実験 1}）。
\item 決算発表文を読ませると、LLM の点数は\textbf{発表をはさむ 2 日間の株価の反応}を強く説明した（1 標準偏差あたり約 +1.2\%）。アナリスト予想との差や、反応で学習した FinBERT を統制しても残った（\textbf{実験 2〜5・10・11}）。
\item しかしその情報は発表日のうちに株価に入りきり、翌日以降や翌月には残っていなかった。月次・日次の BLM は、費用の後で時価総額を有意に上回らなかった（\textbf{実験 7〜9}）。
\end{enumerate}
「LLM は決算発表文を読める。けれども、読めたことは発表の瞬間に市場も知っている」が、この研究の結論である。詳しくは\chapref{ch:qa}と\chapref{ch:conclusion}。
\end{keybox}

\section*{構成}

\begin{longtable}{P{3.8cm}P{11.0cm}}
\toprule
\textbf{部} & \textbf{内容} \\
\midrule
\endhead
\bottomrule
\endfoot
第 I 部　全体の地図 & 問いと答えの要約、2 週間の経緯、すべての実験の\textbf{入力と出力の総覧}。まずここを読む \\
第 II 部　金融工学の基礎 & リターンとリスク、共分散の推定、アクティブ運用（TE・IR・IC）、決算発表とイベントスタディ。新しく書いた章 \\
第 III 部　平均分散と BLM & 平均分散最適化から BLM の公式と主結果まで（解説「平均分散から Black--Litterman へ」と、もう 1 つの解説の導出の章） \\
第 IV 部　ビュー・$\Omega$・$\tau$ & ビューの作り方、確信度、リスクの扱いの要点（解説「BLM のビューとリスクのコツ」） \\
第 V 部　統計・機械学習・言語モデル & 回帰と検定、予測と過学習、文章を数にする方法、LLM の学習と計算、データの時点と先読み。新しく書いた章 \\
第 VI 部　データ基盤 & データベース、取得の仕組み、イベントスタディのデータと点数、追加したデータ \\
第 VII 部　元論文の再現と検証 & 総合報告書 1。元論文の再現、計算の問題、月次レポートと指標、実験 A〜G \\
第 VIII 部　仮説 H01〜H07 & 実験 1。実験の前の設計と、結果と、記録の詳細を仮説ごとに並べた \\
第 IX 部　方針の転換と文献調査 & 総合報告書 2 の前半。Claude・GPT・Gemini の方針と文献調査の比較、計画 \\
第 X 部　決算発表文を読む & 実験 2〜5・10・11。この研究の中心 \\
第 XI 部　ポートフォリオの実験 & 共分散と TE の補足実験（X01〜X05）と、月次・日次の BLM（実験 7〜9） \\
第 XII 部　結論 & 全体の結論、失敗と教訓、残った問い \\
付録 & 費用と計算、用語集、公式、設定、テーブルの定義、プロンプトの全文、コード、事前登録と結果の記録の原文、外部のモデルへの依頼と回答の原文、参考文献 \\
\end{longtable}

\section*{読み方}

\begin{itemize}
\item \textbf{結論だけ知りたい}：\chapref{ch:qa}（問いと答え）、\chapref{ch:io}（入出力の総覧）、\chapref{ch:conclusion}（結論）。
\item \textbf{何をしたかを順に追いたい}：\chapref{ch:history}（経緯）のあと、第 VII 部から第 XI 部を順に。
\item \textbf{背景知識から学びたい}：第 II・III・IV・V 部。各章の最後の「このプロジェクトでの使い方」の囲みが、後の部のどこでその知識を使うかを示す。
\item \textbf{特定の実験を調べたい}：\chapref{ch:io}の表で実験の番号を探し、その章へ。どの実験の章も、見出しのすぐ下に\textbf{入力と出力の囲み}がある。
\end{itemize}

\section*{この本の約束}

\begin{iobox}[入力と出力の囲み（見本）]
\iorow{入力}{モデルに渡したもの（データの種類、期間、件数）}
\iorow{モデル}{使ったモデルの名前と版、設定}
\iorow{学習}{\textbf{何を何に学習させたか}。学習させていない場合はそう書く}
\iorow{出力}{モデルが返したもの}
\iorow{評価}{出力をどう判定したか（目的変数、検定）}
\iorow{結果}{主な数値と判定}
\end{iobox}

\begin{itemize}
\item \textbf{入力と出力の強調}：実験ごとに、上の形の囲みを見出しの下に置いた。LLM などのモデルについては、何を入れて何が出たか、学習させたかどうかを必ず書いた。
\item \textbf{出典}：ほかの文書から取り込んだ章や節には、見出しの下に小さな字で出典を書いた。取り込んだ文章は、元の文書のまま変えていない（見出しの深さと、相互参照の番号だけを変えた）。元の文書の中の「第 N 章」「付録 B」などの参照は、この本の番号に置き換えた。この本に取り込まなかった箇所への参照は「（元の資料の別の箇所）」と書いた。
\item \textbf{日付}：元の文書は、書いた日の時点の知識で書かれている。後の実験で覆った見通しや計画も、そのまま残した（どう覆ったかは、後の章と\chapref{ch:history}にある）。たとえば第 IX 部の「文献調査と設計方針のまとめ」は 10 月 17 日の朝の時点のもので、その後の方針とは違う箇所がある。
\item \textbf{重複}：同じ実験が、総合報告書と個別のレポートの両方に書かれていることがある。この本では原則として詳しいほうを取り、もう一方は省いた。数値が両方にある場合、両者は一致している。
\item \textbf{数値の単位}：bp（ベーシスポイント）は 0.01\%。「1 標準偏差あたり +120bp」は、点数が 1 標準偏差高い発表ほど、反応が平均 1.2\% 高いことを指す。$t$ は $t$ 値、$p$ は $p$ 値。
\item \textbf{秘密の情報}：API の鍵などの秘密の情報は、この本には一切含まれない。データの提供元の利用規約で再配布できない値（アナリスト予想、株価そのものなど）も載せていない。
\end{itemize}

\section*{元になった文書}

\begin{longtable}{P{6.4cm}P{3.0cm}P{5.2cm}}
\toprule
\textbf{文書} & \textbf{書いた日} & \textbf{この本の場所} \\
\midrule
\endhead
\bottomrule
\endfoot
総合報告書 1 & 10/13 & 第 VI・VII 部、付録 \\
解説「平均分散から Black--Litterman へ」 & 10/14 & 第 III 部、第 VII 部、付録 \\
解説「BLM のビューとリスクのコツ」 & 10/14 & 第 III・IV 部、付録 \\
実験仮説（実験 1 の設計書） & 10/15〜16 & 第 VIII 部 \\
実験結果レポート（実験 1） & 10/17 & 第 VIII 部 \\
方針・文献調査の比較とまとめ & 10/17 & 第 IX 部、付録 \\
設計.md（イベントスタディの設計） & 10/17 & 第 VI 部 \\
総合報告書 2 & 10/18 & 第 VI・IX・X 部、付録 \\
実験 2〜4 のレポート & 10/18 & 第 X 部 \\
補足実験 X02〜X04 のレポート、実験 7・8 のレポート & 10/18 & 第 XI 部 \\
実験 5・10 のレポート、実験 9、補足実験 X05 & 10/19 & 第 X・XI 部 \\
設計\_追加データ.md、議論用の案、読ませ方の調査 & 10/19 & 第 VI・X 部 \\
実験 11 の方針と事前登録 & 10/20 & 第 X 部、付録 \\
実験 11 のレポート & 10/21 & 第 X 部 \\
\end{longtable}

この本を作り直すには、\code{資料集/00\_全体の記録} で \code{sh 作成スクリプト/make.sh} を実行する。スクリプト（\code{作成スクリプト/build.py}）は、元の文書から章を取り出し、見出しと参照を直して \code{取り込み/} に置き、新しく書いた章（\code{章/}）と合わせて 1 つの TeX にする。元の文書は変更しない。
